% ================================================================
\section{The Vacuum: $L^2(0,1)$ as a Hilbert Space}
% ================================================================

The proof lives in the Hilbert space $\HH = L^2(0,1)$ with the standard
inner product $\braket{f}{g} = \int_0^1 f(x)\overline{g(x)}\,dx$.

\begin{principle}[The Vacuum State]
The constant function $\mathbf{1}(x) = 1$ on $(0,1)$ plays the role of
the \emph{vacuum state} $\ket{0}$ in the physical theory. The Riemann
Hypothesis is equivalent to the statement that this vacuum can be
\emph{perfectly screened}---approximated to arbitrary precision---by
linear combinations of the basis functions $h_k$.
\end{principle}

In quantum field theory, \emph{vacuum screening} occurs when virtual
particle-antiparticle pairs can completely neutralize a test charge.
In the Cathedral, the ``charge'' is the constant function $1$, and the
``virtual pairs'' are the fractional-part oscillations $\fract{1/(kx)}$
which encode the distribution of primes.

The physical content of RH is therefore:
\begin{quote}
\emph{The prime number vacuum is perfectly screening.}
\end{quote}

\subsection{The Basis Functions as Quantum States}

Each $h_k(x) = \fract{1/(kx)}$ has a sawtooth waveform with $k$ teeth
in the interval $(0,1)$. The Mellin transform
\[
  \mathcal{M}[h_k](s) = \int_0^1 h_k(x) \, x^{s-1}\,dx = \frac{1}{k \cdot (s-1)}
  - \frac{k^{-s}}{s-1} + k^{-s}\mathcal{M}[h_1](s)
\]
reveals a \emph{rank-1 structure}: at any zero $\rho$ of $\zeta$,
\[
  \mathcal{M}[h_k](\rho) = \frac{1}{k(\rho - 1)}.
\]
The $k$-dependence and $\rho$-dependence factorize completely:
$\mathcal{M}[h_k](\rho) = \phi_k \cdot \psi(\rho)$ where
$\phi_k = 1/k$ and $\psi(\rho) = 1/(\rho-1)$.

\begin{correspondence}[Rank-1 = Factorizable Scattering]
The rank-1 factorization $\mathcal{M}[h_k](\rho) = \phi_k \cdot \psi(\rho)$
is the number-theoretic analogue of \emph{factorizable scattering amplitudes}
in integrable quantum field theories. In such theories, the $n$-particle
$S$-matrix decomposes into a product of two-particle amplitudes. Here,
the ``scattering'' of the $k$-th basis function off the zero $\rho$ is
determined entirely by two independent factors: the ``momentum'' $1/k$
and the ``resonance'' $1/(\rho-1)$.
\end{correspondence}

This is proved in the Cathedral as \lean{bd\_mellin\_at\_zero} (BDMellin.lean),
and it is the key lemma in the \emph{Rank-1 Mellin Miracle} that proves
the converse direction of the Nyman--Beurling equivalence.

